\section{The timelike region and the regular static quotient}
\label{geo-section}

We now use the causal stationary field supplied by
Theorem~\ref{eq-causal-killing}.  All geometric objects in this section are
constructed from the original equality data.  The logarithmic estimate below
is a cosmological version of the Killing-norm calculation in
\cite[Section~5.3]{internalflat}; we give its proof, including the treatment of an
arbitrary interior null set.

Write $M=\operatorname{int}\Omega$, and let $\nu$ on $S$ point into $\Omega$.
We record the precise outputs of Theorem~\ref{eq-causal-killing} that are used.
The fields $u,X$ are smooth up to $S$, satisfy $u>0$ on $M$ and $u\geq|X|_g$,
and obey
\begin{equation}
 \operatorname{sym}\nabla X=-uK,
 \qquad X|_S=u\nu,
 \qquad \partial_\nu u+K(X,\nu)=\kappa_h>0,
 \label{geo-input-killing}
\end{equation}
where $\kappa_h$ is a constant.  Either $u|_S>0$ everywhere or $u|_S=0$
everywhere. Writing $\tau'$ for the exponent $\beta$ in
Theorem~\ref{eq-causal-killing}, with $\tau'>3/2$, the end estimates are
\begin{equation}
 u-V_0=O_2(r^{1-\tau'}),\qquad X=O_2(r^{1-\tau'}).
 \label{geo-input-end}
\end{equation}
Here and below the end estimates use the reference hyperbolic tensor norms,
including their stated differentiated H\"older bounds.  The stationary metric
\begin{equation}
 \mathbf g=-u^2\dd t^2+
 g_{ij}(\dd x^i+X^i\dd t)(\dd x^j+X^j\dd t)
 \label{geo-stationary-metric}
\end{equation}
on $\mathbb R\times M$, with $\xi=\partial_t$ and
$N=u^2-|X|_g^2$, satisfies
\begin{equation}
 \Ric_{\mathbf g}=-3\mathbf g+
       8\pi\mu u^{-2}\xi^\flat\otimes\xi^\flat,
 \qquad \mu N=0.
 \label{geo-input-ricci}
\end{equation}
The musical notation on $\xi$ refers to $\mathbf g$; spatial musical notation
will always refer to $g$ unless specified otherwise.

\begin{theorem}[Regular static quotient]
\label{geo-static-quotient}
\label{geo-static-reduction}
Under these hypotheses, $N>0$ throughout $M$, and $M$ is simply connected.
The fields
\begin{equation}
 h=g+N^{-1}X^\flat\otimes X^\flat,
 \qquad \lambda=\sqrt N,\qquad \mathcal A=N^{-1}X^\flat
 \label{geo-quotient-fields}
\end{equation}
satisfy $h\geq g$ and $\dd\mathcal A=0$, and there is a global smooth
function $H$ on $M$
with $\dd H=-\mathcal A$.  They obey
\begin{equation}
 \Delta_h\lambda=3\lambda,
 \qquad \Hess_h\lambda=\lambda(\Ric_h+3h),
 \qquad R_h=-6.
 \label{geo-static-equations}
\end{equation}
Attaching $S$ in the regular boundary coordinates of
Lemma~\ref{geo-boundary-collars} gives a smooth complete Riemannian manifold
with compact connected nondegenerate boundary.  Its underlying completion is homeomorphic to
$\Omega$, and its boundary is identified diffeomorphically with $S$.  On that
boundary,
\begin{equation}
 h|_{TS}=g|_{TS},\qquad \Area_h(S)=\Area_g(S),\qquad \lambda=0,
 \qquad \partial_{\nu_h}\lambda=\kappa_h,
 \qquad \mathrm{II}_h=0.
 \label{geo-quotient-boundary}
\end{equation}
Moreover, in the original end chart,
\begin{equation}
 h-g=O_2(r^{-2\tau'}),\qquad
 \lambda-V_0=O_2(r^{1-\tau'}),\qquad \frac{\lambda}{V_0}\longrightarrow1.
 \label{geo-quotient-end}
\end{equation}
\end{theorem}

\subsection{Boundary collars and the Killing identities}

Initially put $P=\{N>0\}\subset M$.  On $P$ define $h,\lambda,\mathcal A$ as
in Equation~\eqref{geo-quotient-fields}, and set
\begin{equation}
 C=h^{-1}=g^{-1}-u^{-2}X\otimes X,
 \qquad \mathcal F=\dd\mathcal A.
 \label{geo-C-F}
\end{equation}
The tensor $C$ extends smoothly and nonnegatively to all of $M$.  Direct
completion of the square gives
\begin{equation}
 \mathbf g=h-\lambda^2(\dd t-\mathcal A)^2,
 \qquad \dd V_h=\frac{u}{\lambda}\dd V_g.
 \label{geo-square-completion}
\end{equation}

\begin{lemma}[The end and the original boundary collar]
\label{geo-positive-collars}
The end and a collar of $S$ in $M$ lie in $P$.  In particular the interior
zero set $Z=\{N=0\}\cap M$ is compact.  If $u|_S>0$, then
\begin{equation}
 \dd N|_S=2\kappa_h X^\flat|_S,
 \qquad N=2\kappa_h u|_S\,s+O(s^2),
 \label{geo-positive-boundary-jet}
\end{equation}
where $s\geq0$ is $g$-normal distance into the domain.  If $u|_S=0$, there
are smooth $a$ and $W$ in this collar such that
\begin{equation}
 u=sa,\qquad a|_S=\kappa_h,\qquad X=s^2W,
 \qquad N=s^2\bigl(a^2-s^2|W|_g^2\bigr).
 \label{geo-zero-boundary-jet}
\end{equation}
\end{lemma}
\begin{proof}
Since $X=u\nu$ on $S$, the function $N$ vanishes there.  The normal component
of Equation~\eqref{geo-input-killing} gives
$\langle\nabla_\nu X,\nu\rangle=-uK(\nu,\nu)$, and hence
\[
 \partial_\nu N
 =2u\partial_\nu u-2u\langle\nabla_\nu X,\nu\rangle
 =2u\bigl(\partial_\nu u+K(X,\nu)\bigr)=2u\kappa_h.
\]
All tangential derivatives of $N$ vanish on $S$.  This proves
Equation~\eqref{geo-positive-boundary-jet} and the positive-lapse collar
assertion.  If $u=0$ on $S$, then $X=0$ there and its tangential covariant
derivatives vanish.  The normal-normal and normal-tangential components of
$\operatorname{sym}\nabla X=0$ give $\nabla_\nu X=0$ on $S$.  Smooth division
by $s$ and by $s^2$ gives Equation~\eqref{geo-zero-boundary-jet}; the boundary
lapse equation gives $a|_S=\kappa_h$.  Compactness makes the collars uniform.
Finally, Equation~\eqref{geo-input-end} gives
$N=V_0^2(1+O_2(r^{-\tau'}))$, so the entire sufficiently distant end lies in
$P$.  The remaining zero set is a closed subset of a compact set separated
from $S$.
\end{proof}

\begin{lemma}[Norm and twist identities]
\label{geo-killing-identities}
Let $D_{\alpha\beta}=\nabla^{\mathbf g}_\alpha\xi_\beta$, with the full
Lorentzian contraction used in $|D|_{\mathbf g}^2$.  On $M$,
\begin{equation}
 \frac1u\operatorname{div}_g(uC\,\dd N)
       =6N-2|D|_{\mathbf g}^2.
 \label{geo-norm-identity}
\end{equation}
On $P$, the antisymmetric contravariant tensor
\begin{equation}
 Q^{ij}=uN C^{ik}C^{j\ell}\mathcal F_{k\ell}
 \label{geo-Q}
\end{equation}
satisfies $\nabla_iQ^{ij}=0$.
\end{lemma}
\begin{proof}
The Killing equation makes $D$ antisymmetric.  Killing differentiation gives
$\Box_{\mathbf g}\xi_\beta=-\Ric_{\mathbf g\,\beta\gamma}\xi^\gamma$.
Since $\mathbf g(\xi,\xi)=-N$, it follows that
\[
 \Box_{\mathbf g}N
 =-2|D|_{\mathbf g}^2+2\Ric_{\mathbf g}(\xi,\xi)
 =-2|D|_{\mathbf g}^2+6N.
\]
In the last equality the additional term in
Equation~\eqref{geo-input-ricci} vanishes by $\mu N=0$.  For a stationary
scalar the spatial block of $\mathbf g^{-1}$ is $C$, and the spacetime volume
density is $u\sqrt{\det g}$.  This proves Equation~\eqref{geo-norm-identity}.

For the second identity, write $\vartheta=\dd t-\mathcal A$.
The Killing one-form is $\xi^\flat=-\lambda^2\vartheta$, so
$\xi^\flat\wedge\dd\xi^\flat=-\lambda^4\vartheta\wedge\mathcal F$.
The differentiated Killing equation, contracted with the volume form, gives
\[
 \dd *_{\mathbf g}(\xi^\flat\wedge\dd\xi^\flat)
       =2*_{\mathbf g}\bigl(\xi^\flat\wedge\Ric_{\mathbf g}(\xi,\cdot)\bigr)
       =0.
\]
The right side vanishes because $\Ric_{\mathbf g}(\xi,\cdot)$ is proportional
to $\xi^\flat$.  Evaluating the Hodge star using the unit time form
$\lambda\vartheta$ turns the left side, up to a fixed orientation sign, into
$\dd(\lambda^3 *_h\mathcal F)$.  Thus
\[
 \operatorname{div}_h\bigl(\lambda^3 C^{ik}C^{j\ell}
                                      \mathcal F_{k\ell}\bigr)=0.
\]
Use $\lambda^3\dd V_h=uN\dd V_g$ to change the divergence density.  The
connection term on the free index of an antisymmetric two-tensor vanishes
for either torsion-free connection.  The resulting equation is exactly
$\nabla_iQ^{ij}=0$.
\end{proof}

\subsection{Logarithmic control and vanishing twist}

\begin{lemma}[Transverse logarithmic energy]
\label{geo-logarithmic-energy}
On a neighborhood of $Z$, set $B=|X|_g$, $e=X/B$, and $T=e^\perp$.
For each smooth compactly supported cutoff $\eta$ in that neighborhood,
\begin{equation}
 \sup_{\varepsilon>0}
 \int_M u\eta^2\frac{|\dd N|_C^2}{(N+\varepsilon)^2}\dd V_g<\infty.
 \label{geo-log-energy-bound}
\end{equation}
In particular, $\dd_T\log N$ is locally square integrable on the positive
side up to $Z$:
\begin{equation}
 \int_P\eta^2|\dd_T\log N|_g^2\dd V_g<\infty.
 \label{geo-log-transverse}
\end{equation}
No regularity or measure assumption on $Z$ is required.
\end{lemma}
\begin{proof}
On a sufficiently small fixed neighborhood of the compact set $Z$, both
$u$ and $B$ are bounded below by positive constants.  The tensor $C$ acts as
the identity on $T$ and has eigenvalue $N/u^2$ in direction $e$.
Smooth nonnegativity of $N$ implies
\begin{equation}
 |\dd N|_g^2\leq C_0N
 \label{geo-quadratic-vanishing}
\end{equation}
on compact subsets after slightly enlarging the neighborhood.  Indeed,
choose a uniform coordinate radius and a bound $L$ for $\Hess_gN$ large
enough that $|\dd N|/L$ is less than that radius.  Taylor's Inequality along
the negative gradient geodesic at length $|\dd N|/L$ yields
$0\leq N-|\dd N|^2/(2L)$.

Put $n=(\partial_t-X)/u$, and complete $e$ to a spatial orthonormal frame
$e,E_1,E_2$.  Since $\xi=un+Be$, differentiation of its norm gives
\[
 uD_{in}=-\tfrac12N_i-BD_{ie}.
\]
Writing $v_a=D_{ae}$, expansion of the temporal and spatial contractions
therefore gives the exact formula
\begin{equation}
 -2|D|_{\mathbf g}^2
 =\frac{|\dd N|_g^2}{u^2}
   +\frac{4B}{u^2}\sum_{a=1}^2N_av_a
   -\frac{4N}{u^2}\sum_{a=1}^2v_a^2-4D_{12}^2.
 \label{geo-D-expansion}
\end{equation}
All coefficients and the components of $D$ are bounded on the fixed
neighborhood.  Equations~\eqref{geo-norm-identity} and
\eqref{geo-quadratic-vanishing} consequently imply
\begin{equation}
 \frac1u\operatorname{div}_g(uC\,\dd N)
                  \leq C_1(N+|\dd_TN|_g).
 \label{geo-norm-upper-bound}
\end{equation}
The additional cosmological term $6N$ has the harmless first order shown.

Multiply Equation~\eqref{geo-norm-upper-bound} by
$u\eta^2/(N+\varepsilon)$ and integrate. Writing $E_\varepsilon$ for the
integral in Equation~\eqref{geo-log-energy-bound}, before taking the
supremum, integration by parts gives
\[
 E_\varepsilon
 -2\int_M\frac{u\eta C(\dd\eta,\dd N)}{N+\varepsilon}\dd V_g
 \leq C_1\int_Mu\eta^2\dd V_g
    +C_1\int_M\frac{u\eta^2|\dd_TN|}{N+\varepsilon}\dd V_g.
\]
Cauchy's Inequality bounds the cutoff term by
$E_\varepsilon/4+C\int u|\dd\eta|_C^2$, and the last term by
$E_\varepsilon/4+C\int u\eta^2$, because
$|\dd_TN|\leq|\dd N|_C$.  This proves the uniform bound.  Fatou's Lemma on
$P$, together with the positive lower bound for $u$, proves
Equation~\eqref{geo-log-transverse}.
\end{proof}

\begin{lemma}[Vanishing twist]
\label{geo-vanishing-twist}
The form $\mathcal A$ is closed on all of $P$.
\end{lemma}
\begin{proof}
Lower the free index of $Q^{kj}\mathcal A_j$ using $g$.  Direct substitution
of
\[
 \mathcal F=N^{-1}\dd X^\flat-N^{-2}\dd N\wedge X^\flat,
 \qquad C X^\flat=(N/u^2)X
\]
gives
\begin{equation}
 W_i:=g_{ik}Q^{kj}\mathcal A_j
  =u^{-1}\left\{X^j(\dd X^\flat)_{ij}
                    -B^2(\dd_T\log N)_i\right\}.
 \label{geo-twist-contraction}
\end{equation}
This expression is orthogonal to $X$.  At a zero of $X$ its second numerator
is interpreted as $B^2\dd N-\dd N(X)X^\flat$, so the identity does not
require a choice of $e$ away from the neighborhood used above.

Let $\chi$ be a product of three nonnegative cutoffs, excluding first $Z$,
then the original boundary, then infinity.  Testing $\nabla_iQ^{ij}=0$ by
$\chi\mathcal A_j$ gives
\begin{equation}
 \frac12\int_P\chi uN C^{ik}C^{j\ell}
                 \mathcal F_{ij}\mathcal F_{k\ell}\dd V_g
       =-\int_P\langle W,\dd\chi\rangle_g\dd V_g.
 \label{geo-twist-energy}
\end{equation}
The integrand on the left is nonnegative.

Keep the boundary and infinity cutoffs fixed.  Let the first cutoff change
from zero to one on $\varepsilon<N<2\varepsilon$, with derivative bounded
by $C/\varepsilon$.  Orthogonality in
Equation~\eqref{geo-twist-contraction} makes its error bounded by
\[
 C\int_{\{\varepsilon<N<2\varepsilon\}}
       \bigl(|\dd_T\log N|+|\dd_T\log N|^2\bigr)\dd V_g
\]
on a fixed compact set.  This tends to zero by
Lemma~\ref{geo-logarithmic-energy}: the shell indicators tend pointwise to
zero on $P$ and are dominated by an integrable function there.  The argument
also applies if $Z$ has positive volume.

For the boundary cutoff use the normal coordinate $s$.  If $u|_S>0$,
Equation~\eqref{geo-positive-boundary-jet} gives $N\asymp s$,
$X=u|_S\nu+O(s)$, and $\dd_TN=O(s)$.  Thus $W$ is bounded.  Its
orthogonality to $X$ implies $W(\nu)=O(s)$, so a cutoff with derivative
$O(\delta^{-1})$ on $\delta<s<2\delta$ has error $O(\delta)$.  If $u|_S=0$,
Equation~\eqref{geo-zero-boundary-jet} gives
\[
 u\asymp s,\quad X=O(s^2),\quad \dd X^\flat=O(s),
 \quad N\asymp s^2,\quad \dd N=O(s),
\]
and Equation~\eqref{geo-twist-contraction} gives $W=O(s^2)$, with the same
conclusion.

Finally, Equation~\eqref{geo-input-end} gives
$W=O(r^{1-2\tau'})$ on the end.  A cutoff on $R<r<2R$ has bounded
hyperbolic gradient and the shell has volume $O(R^2)$.  Its error is
$O(R^{3-2\tau'})$, which vanishes since $\tau'>3/2$.  Successively removing
these cutoffs and using Fatou's Lemma in
Equation~\eqref{geo-twist-energy} makes its nonnegative integrand vanish
everywhere.  Since $uN>0$ and $C$ is positive definite on $P$, one has
$\mathcal F=0$ there.
\end{proof}

\begin{lemma}[Static equations and connectedness]
\label{geo-static-on-P}
Equations~\eqref{geo-static-equations} hold on $P$, and $P$ is connected.
\end{lemma}
\begin{proof}
Locally write $\mathcal A=\dd f$ and $T_0=t-f$.  Then
$\mathbf g=h-\lambda^2\dd T_0^2$.  Equation~\eqref{geo-input-ricci} is
vacuum on $P$.  The only mixed-time Christoffel symbols of this static
product are
\[
 \Gamma^{T_0}_{iT_0}=\lambda^{-1}\partial_i\lambda,
 \qquad \Gamma^i_{T_0T_0}=\lambda h^{ij}\partial_j\lambda.
\]
Thus its spatial Ricci tensor is
$\Ric_h-\lambda^{-1}\Hess_h\lambda$, and its time-time component is
$\lambda\Delta_h\lambda$.  Comparing with $-3\mathbf g$ gives the first
two equations in Equation~\eqref{geo-static-equations}; tracing gives
$R_h=-6$.  These equations are intrinsic and hence hold globally on $P$.

The end belongs to a single component of $P$.  Any other component has
compact closure in the original manifold with boundary.  The continuous
function $\lambda$ is zero on every finite boundary point of that component,
including a possible point on $S$.  It therefore attains a positive maximum
at an interior point, contrary to $\Delta_h\lambda=3\lambda$.  There is no
other component.
\end{proof}

\subsection{Optical completeness and simple connectivity}

\begin{lemma}[Completeness after truncation of the ordinary end]
\label{geo-optical-completeness}
Equip $P$ with the optical metric $k=\lambda^{-2}h$.  All escape directions
other than the ordinary end have infinite $k$-length.  After truncating the
ordinary end at a sufficiently large sphere, the resulting manifold with
boundary is complete for its intrinsic optical distance.  The same is true
of each of its Riemannian covering spaces.
\end{lemma}
\begin{proof}
Near $Z$, smooth nonnegativity and compactness give
$N\leq C\operatorname{dist}_g(\cdot,Z)^2$.  If
$\rho=\operatorname{dist}_g(\cdot,Z)$, then $|\dd\rho|_g\leq1$ almost
everywhere and
\[
 |\dd\log\rho|_k^2
   =N C(\dd\log\rho,\dd\log\rho)
   \leq N\rho^{-2}\leq C.
\]
Hence approach to $Z$ has infinite optical length.  In the zero boundary
lapse case the same calculation uses $s$ and $N=O(s^2)$.

For positive boundary lapse it is essential to allow arbitrary tangential
motion.  In the $g$-normal collar decompose
$X=X^\nu\nu+X_T$.  Since $|\dd s|_g=1$, the exact inverse-metric formula
is
\begin{equation}
 |\dd s|_k^2
 =N\left(1-\frac{(X^\nu)^2}{u^2}\right)
 =\frac{N}{u^2}\bigl(N+|X_T|_g^2\bigr)=O(s^2).
 \label{geo-optical-arbitrary-curve}
\end{equation}
Here $N\asymp s$, $X_T=O(s)$, and $u$ is bounded below.  Therefore
$|\dd s(v)|\leq Cs|v|_k$ for every vector $v$, and any curve approaching
$S$ has infinite optical length.  This proves the assertion without
restricting to normal curves.

After truncation, $N$ is bounded above, so $k\geq c g$.  Any optical Cauchy
sequence is consequently Cauchy for $g$ and has a limit in the compact
original region with the large sphere attached.  The logarithmic estimates
exclude a limit in $Z$ or in $S$.  Every other limit is an ordinary smooth
point of the truncated metric, including its large-sphere boundary.  This
proves completeness.  A smooth manifold with boundary and its intrinsic
Riemannian distance is a locally compact length space; completeness makes
it proper.  For a covering metric completeness follows by projecting a
Cauchy sequence and then taking an evenly covered small neighborhood of its
limit.  Sufficiently short connecting paths remain in that neighborhood, so
the sequence eventually lies in a single sheet and converges there.
\end{proof}

\begin{lemma}[Two convex faces cannot be joined by a shortest optical path]
\label{geo-null-index-obstruction}
Let $(L,h,\lambda)$ satisfy the static equations in
Equation~\eqref{geo-static-equations}, with $\lambda>0$.  Suppose that two
compact boundary faces $\Sigma_0,\Sigma_1$ have positive $h$-mean curvature
for the normals pointing out of $L$.  There is no shortest optical path
between these faces which is orthogonal at its endpoints and otherwise
interior.
\end{lemma}
\begin{proof}
Suppose such a path has optical length $\ell>0$.  In the abstract static
product $(\mathbb R\times L,\widehat{\mathbf g})$, where
$\widehat{\mathbf g}=-\lambda^2\dd T^2+h$, its lift with increasing $T$
from $0$ to $\ell$ is a null geodesic up to reparametrization.  One can see
this directly from the null energy relation
$\lambda^2\dot T=E>0$: optical arclength and affine parameter are related by
$\dd s/\dd T=\lambda^2/E$, and the spatial geodesic equation becomes the
optical geodesic equation.  Use an affine parameter on a compact interval
$[a,b]$ and write its tangent as $k_0$.

This null curve has no future timelike competitor with endpoints in the
fixed-time surfaces
$\{0\}\times\Sigma_0$ and $\{\ell\}\times\Sigma_1$.  Indeed every such
competitor satisfies
\[
 \ell=\int_a^b\dot T\,\dd s
 >\int_a^b\frac{|\dot x|_h}{\lambda}\,\dd s
 \geq\operatorname{dist}_{\lambda^{-2}h}(\Sigma_0,\Sigma_1)=\ell.
\]
We produce a competitor by second variation.

Choose a parallel orthonormal screen pair $E_1,E_2$ along the null geodesic,
with $\langle E_A,k_0\rangle=0$, initially tangent to
$\{0\}\times\Sigma_0$.  At the terminal endpoint, add suitable multiples
of $k_0$ to remove their time components.  Orthogonality to $k_0$ then
places them in $T\Sigma_1$.  Extending these multiples smoothly along the
curve gives fields $V_A=E_A+b_Ak_0$ tangent to the endpoint surfaces, still
orthonormal, with
$\langle\nabla_{k_0}V_A,\nabla_{k_0}V_A\rangle=0$.
Let $H_0,H_1>0$ denote the specified mean curvatures.  The spatial component
of $k_0$ is $-c_0\nu_0$ initially and $c_1\nu_1$ terminally, for positive
$c_0,c_1$.  The static time slices have zero second fundamental form.  The
endpoint terms in the energy second variation therefore have trace
$-c_0H_0-c_1H_1$.

More explicitly, choose variations with initial and terminal curves in the
fixed-time surfaces and first variational fields $V_A$.  The second
variation of $\frac12\int\langle\dot\gamma,\dot\gamma\rangle$ is
\[
 I(V_A,V_A)=
 \int_a^b\left(
 |\nabla_{k_0}V_A|^2
 -\langle R(V_A,k_0)k_0,V_A\rangle\right)\dd s
 +\left[\langle\nabla_{V_A}V_A,k_0\rangle\right]_a^b.
\]
At an endpoint, the surface acceleration has normal component
$-\mathrm{II}_{\nu_j}(V_A,V_A)$.  This proves the stated signs.  Generator
multiples do not change the curvature term, and summing the screen
curvatures gives $\Ric_{\widehat{\mathbf g}}(k_0,k_0)=0$.  Consequently
\begin{equation}
 \sum_{A=1}^2I(V_A,V_A)=-c_0H_0-c_1H_1<0.
 \label{geo-negative-null-index}
\end{equation}
Choose one of these fields with negative index and a smooth variation
$\gamma_\varepsilon$ realizing it.  The first variation of squared speed
vanishes pointwise, because
$\langle\nabla_{k_0}V,k_0\rangle=0$.  Put
$q(s)=\partial_\varepsilon^2|\dot\gamma_\varepsilon|^2|_{\varepsilon=0}$;
its integral is $2I(V,V)<0$.  Let $\overline q$ be its negative average.
Choose a parallel auxiliary null vector $L_0$ with
$\langle L_0,k_0\rangle=-1$, and add to the variation a second-order
displacement $\frac12\varepsilon^2 f(s)L_0$, where
\[
 f(a)=0,\qquad f'(s)=\tfrac12(q(s)-\overline q).
\]
Then $f(b)=0$, so the endpoints are unchanged; the second variation of
squared speed changes by $-2f'$ and becomes the constant $\overline q<0$.
Uniform Taylor expansion on $[a,b]$ now makes the varied curve timelike for
all sufficiently small nonzero $\varepsilon$.  It remains future directed
because $\dot T>0$ on the original compact segment.  Its endpoint contacts
remain transverse and directed into and out of $L$; away from small endpoint
intervals the original curve lies a positive distance from the boundary.
Thus the variation remains in $L$ and is the forbidden timelike competitor.
\end{proof}

\begin{proposition}[Simple connectivity of the timelike region]
\label{geo-simple-connectivity}
The connected manifold $P$ is simply connected.
\end{proposition}
\begin{proof}
Suppose its universal cover is nontrivial.  The ordinary end is a product of
a sphere and a ray and is simply connected.  Its inverse image therefore
consists of at least two disjoint copies.  Truncate every copy at the same
large radius $R$.  Their boundary faces are compact spheres.  By
Equation~\eqref{geo-input-end} and the formula for $h$, their mean curvatures
for the normals pointing into the removed tails satisfy
\[
 H_h=\frac{2\sqrt{1+R^2}}{R}+o(1)>0.
\]
The truncated cover is connected: attaching or removing a product tail
cannot join different components of its complement.  By
Lemma~\ref{geo-optical-completeness} its intrinsic optical metric is complete
and proper.

Fix one face.  The union of all the other faces is a closed set, since it
is the inverse image of the truncation sphere with the fixed component
removed: local finiteness in evenly covered neighborhoods makes each
component open as well as closed within that inverse image.  The remaining
union is disjoint from the fixed compact face.  Properness implies that the positive
distance between that face and the union is attained.  A minimizing path
has no intermediate contact with a face: a later contact with the initial
face or an earlier contact with any other face would shorten the selected
distance. To check its endpoint directions without assuming optical
convexity, let \(\rho\) be the inward \(k\)-distance to the initial
face in its smooth normal collar, and parametrize the path by unit
optical speed from \(t=0\). The bound \(|d\rho|_k=1\) gives
\(\rho(\gamma(t))\le t\); replacing the initial segment by a normal
collar segment gives the reverse inequality by minimality. Hence
\(\rho(\gamma(t))=t\) and \(\dot\gamma=\nabla_k\rho\) near the
endpoint. Apply the same argument backwards at the final face.
The path is therefore an interior geodesic with transverse orthogonal
endpoint contacts. The static equations lift to the cover, so
Lemma~\ref{geo-null-index-obstruction} gives a contradiction.
\end{proof}

\subsection{The geometry of a possible interior null set}

\begin{lemma}[Foliated coordinates and boundary-leaf equations]
\label{geo-null-leaf-equations}
There is a neighborhood of $Z$ on which $X^\flat$ defines a smooth
cooriented integrable plane field.  In local foliated coordinates it has the
form
\begin{equation}
 X^\flat=B_1\dd y,\qquad B_1>0,\qquad
 g=q_y+D_1^2\dd y^2,\qquad N=B_1a(y),\qquad a\geq0.
 \label{geo-null-foliation}
\end{equation}
At any leaf in the relative boundary of $P$, one has $a=a'=0$.  Put
$b_1=\sqrt{B_1}$ and let $\mathcal K_y$ denote the Gaussian curvature of
$q_y$.  On such a boundary leaf,
\begin{equation}
 \Hess_{q_y}b_1=\tfrac12b_1(\mathcal K_y+3)q_y,
 \qquad \Delta_{q_y}(b_1^2)=6b_1^2-a''.
 \label{geo-boundary-leaf-system}
\end{equation}
\end{lemma}
\begin{proof}
The compactness of $Z$ and positivity of $u$ give $X\ne0$ on a neighborhood
of $Z$.  On $P$, closedness of $\mathcal A$ gives
$X^\flat\wedge\dd X^\flat=0$.  This extends by continuity to the relative
boundary of $P$.  In an open region where $N=0$, Equation~\eqref{geo-D-expansion}
and Equation~\eqref{geo-norm-identity} give $D_{12}=0$.  Since the spatial
restriction of $\dd\xi^\flat$ is $\dd X^\flat=2D_{ij}\dd x^i\otimes\dd x^j$
in components, the restriction of $\dd X^\flat$ to $X^\perp$ vanishes there
also.  These cases exhaust the neighborhood, proving Frobenius integrability.

Choose the transverse coordinate oriented so that $X^\flat=B_1\dd y$ with
$B_1>0$, and transport coordinates within the leaves orthogonally using $g$.
This gives the metric decomposition in Equation~\eqref{geo-null-foliation}.
On the positive set, $\dd(B_1\dd y/N)=0$ says that every leaf derivative of
$N/B_1$ vanishes.  The same derivatives vanish on the interior of the zero
set and, by continuity, on its remaining boundary.  Thus $N/B_1=a(y)$
throughout the coordinate neighborhood.  Nonnegativity gives $a=a'=0$ at
a zero.  A boundary leaf is approached by leaves with $a>0$.

We compute the limit of the static equations from those positive leaves.
On them,
\[
 h=q_y+P_1^2\dd y^2,\qquad
 P_1=\left(D_1^2+\frac{B_1}{a}\right)^{1/2},
 \qquad \lambda=b_1\sqrt a.
\]
The second fundamental form of a leaf in $h$ is
$k_{AB}=(2P_1)^{-1}\partial_yq_{AB}=O(\sqrt a)$, and
$P_1^{-1}\partial_yk_{AB}=O(a+|a'|)$.  The tangential Ricci formula is
therefore
\[
 (\Ric_h)_{AB}
  =\mathcal K_y q_{AB}-P_1^{-1}(\Hess_{q_y}P_1)_{AB}
        +O(a+|a'|)
  \longrightarrow
  \mathcal K_y q_{AB}-b_1^{-1}(\Hess_{q_y}b_1)_{AB}.
\]
The tangential normalized Hessian has the limit
\[
 \lambda^{-1}(\Hess_h\lambda)_{AB}
 =b_1^{-1}(\Hess_{q_y}b_1)_{AB}+O(a+|a'|)
 \longrightarrow b_1^{-1}(\Hess_{q_y}b_1)_{AB}.
\]
These estimates follow directly by inserting $P_1=b_1a^{-1/2}
(1+aD_1^2/B_1)^{1/2}$; all coefficients in their error terms are bounded
smooth functions on a fixed coordinate patch.  The tangential static
Ricci equation gives the first equation in
Equation~\eqref{geo-boundary-leaf-system}.

For the scalar equation let $v=\sqrt{\det q_y}$, $f=\log b_1$, and
$w=(1+aD_1^2/B_1)^{1/2}$.  The normal divergence contribution to
$\lambda^{-1}\Delta_h\lambda$ is exactly
\begin{equation}
 \frac1{b_1^2wv}\partial_y
       \left[vw^{-1}\left(af_y+\frac{a'}2\right)\right].
 \label{geo-normal-laplacian-limit}
\end{equation}
At $a=a'=0$ it tends to $a''/(2b_1^2)$.  The tangential contribution tends
to $\Delta_{q_y}b_1/b_1+|\dd b_1|_{q_y}^2/b_1^2$.  Hence
$\Delta_h\lambda=3\lambda$ gives
\[
 b_1\Delta_{q_y}b_1+|\dd b_1|_{q_y}^2+\frac{a''}2=3b_1^2,
\]
which is the second equation in Equation~\eqref{geo-boundary-leaf-system}.
\end{proof}

\subsection{Exclusion of flat boundary null leaves}
\label{geo-flat-subsection}

Here a boundary null leaf is called flat when $\Hess_gN=0$ along it,
equivalently when $a''(y_0)=0$ in Equation~\eqref{geo-null-foliation}.
We exclude these leaves first; the remaining null leaves are then ruled
out in Proposition~\ref{geo-no-null-set}.

We first give the elementary orbitwise construction needed at a possibly
singular compact invariant set.  The construction produces an arc through
each point separately and requires no manifold structure on that set.

\begin{lemma}[A transverse arc with fast backward decay]
\label{geo-flat-fast-arc}
Let $\phi_t$ be a smooth complete flow on a smooth Riemannian
three-manifold $(M,g)$, and let $\mathcal K\subset M$ be a compact set
invariant for all $t\in\mathbb R$.  Suppose that a smooth nonvanishing
one-form $\alpha$ is defined near $\mathcal K$, that
$E=\ker\alpha$ is preserved by the differential of the flow along
$\mathcal K$, and that, for every $p\in\mathcal K$, $t\geq0$, and
$v\in E_p$,
\begin{equation}
\label{geo-flat-rate-hypotheses}
 \lvert v\rvert_g\leq
 \lvert D\phi_t(p)v\rvert_g\leq e^t\lvert v\rvert_g,
 \qquad
 \alpha_{\phi_t(p)}\circ D\phi_t(p)=e^{2t}\alpha_p.
\end{equation}
There are constants $\epsilon,C>0$ such that every $p\in\mathcal K$
admits a $C^1$ embedded arc
$\gamma_p:(-\epsilon,\epsilon)\longrightarrow M$ with
\begin{equation}
\label{geo-flat-arc-conclusion}
 \gamma_p(0)=p,\qquad
 \alpha_p(\gamma_p'(0))=1,\qquad
 d_g\bigl(\phi_{-t}(\gamma_p(\xi)),\phi_{-t}(p)\bigr)
 \leq C\lvert\xi\rvert e^{-3t/2}\quad(t\geq0).
\end{equation}
By decreasing $\epsilon$, all the backward arcs can be kept in any
prescribed neighborhood of $\mathcal K$.  No regularity of the family
$p\mapsto\gamma_p$ is asserted or needed.
\end{lemma}

\begin{proof}
\emph{An invariant transverse line.}
Write $f=\phi_1$ and use the smooth complementary vector
\[
 \nu_p=\frac{\alpha_p^\sharp}{\lvert\alpha_p\rvert_g^2},
 \qquad \alpha_p(\nu_p)=1.
\]
Relative to $T_pM=E_p\oplus\mathbb R\nu_p$, the forward derivative
has the block form
\begin{equation}
\label{geo-flat-forward-blocks}
 Df_p=
 \begin{pmatrix} A_p&B_p\\0&e^2\end{pmatrix},
 \qquad \lVert A_p\rVert\leq e.
\end{equation}
Here $B_p\in E_{f(p)}$ and $\sup_{\mathcal K}\lvert B_p\rvert_g<\infty$.
Fix a full orbit $p_n=f^n(p)$, $n\in\mathbb Z$.  A transverse line at
$p_n$ with generator $\nu_{p_n}+v_n$, $v_n\in E_{p_n}$, is invariant
precisely when its slopes satisfy
\begin{equation}
\label{geo-flat-slope-recurrence}
 v_{n+1}=e^{-2}(A_{p_n}v_n+B_{p_n}).
\end{equation}
Each affine map in this recurrence contracts differences by at most
$e^{-1}$.  Starting with slope zero at $p_{n-k}$ and letting
$k\to\infty$ therefore gives a bounded solution at $p_n$; the solutions
at different indices obey the same recurrence.  Equivalently, one may
sum the absolutely convergent series obtained by repeated substitution.
Its bound is
\begin{equation}
\label{geo-flat-slope-bound}
 \sup_{n\in\mathbb Z}\lvert v_n\rvert_g
 \leq\frac{\sup_{\mathcal K}\lvert B\rvert_g}{e^2-e}.
\end{equation}
The bounded solution is unique: the difference of two such solutions at
index $n$ is at most $e^{-k}$ times their uniformly bounded difference
at index $n-k$.

Set $u_n=\nu_{p_n}+v_n$ and $U_n=\mathbb R u_n$.  Then
\[
 \alpha_{p_n}(u_n)=1,\qquad Df_{p_n}u_n=e^2u_{n+1}.
\]
Compactness, nonvanishing of $\alpha$, and
Equation~\eqref{geo-flat-slope-bound} give uniform equivalence between
the ambient norm and the norm
\[
 \lvert(s,w)\rvert_n=\max\{\lvert s\rvert,\lvert w\rvert_g\},
 \qquad su_n+w\in U_n\oplus E_{p_n}.
\]
In particular these transverse lines have angles uniformly bounded away
from zero relative to $E$.

\emph{Uniform charts and backward products.}
Put $q_j=p_{-j}$ for $j\geq0$.  Choose centered charts
\[
 \Psi_j(s,w)=\exp_{q_j}(s u_{-j}+w),\qquad w\in E_{q_j}.
\]
They and their inverses have uniformly bounded derivatives on suitable
fixed radii.  This follows from compactness and the uniformly bounded
linear coordinate changes just constructed.  No derivatives of $U_n$
as a function of $p_n$ are used.  On a fixed smaller radius $\rho>0$,
the backward time-one maps in these charts have the form
\begin{equation}
\label{geo-flat-backward-maps}
 \Psi_{j+1}^{-1}\circ\phi_{-1}\circ\Psi_j(z)
 =\begin{pmatrix}c&0\\0&T_j\end{pmatrix}z+R_j(z),
 \qquad c=e^{-2},
\end{equation}
where $T_j=D\phi_{-1}|_{E_{q_j}}$ maps $E_{q_j}$ onto $E_{q_{j+1}}$.
The remainders satisfy, in the displayed norms,
\begin{equation}
\label{geo-flat-remainders}
 R_j(0)=0,\qquad DR_j(0)=0,\qquad
 \sup_{j\geq0,\,\lvert z\rvert\leq\rho}
 \lVert D^2R_j(z)\rVert\leq M_0
\end{equation}
for some finite $M_0$.  We increase $M_0$ to a positive number if
necessary.  Equation~\eqref{geo-flat-rate-hypotheses} gives
\begin{equation}
\label{geo-flat-inverse-products}
 \bigl\lVert T_j^{-1}T_{j+1}^{-1}\cdots T_i^{-1}\bigr\rVert
 \leq e^{i-j+1}\qquad(i\geq j).
\end{equation}
Indeed that product is the forward leaf derivative from $q_{i+1}$ to
$q_j$.

\emph{The two Green sums.}
Fix $\beta=3/2$ and consider the Banach space of sequences
$z_j=(s_j,w_j)\in\mathbb R\oplus E_{q_j}$ with norm
\[
 \lVert z\rVert_\beta
 =\sup_{j\geq0}e^{\beta j}\max\{\lvert s_j\rvert,\lvert w_j\rvert_g\}.
\]
For a forcing sequence $r_j=(r_j^s,r_j^E)$ taking values in
$\mathbb R\oplus E_{q_{j+1}}$, use the analogous norm with weight
$e^{\beta j}$.  Define
\begin{align}
\label{geo-flat-green-sums}
 (Gr)^s_j
 &=\sum_{i=0}^{j-1}c^{j-1-i}r_i^s,\\
 (Gr)^E_j
 &=-\sum_{i=j}^{\infty}
 T_j^{-1}T_{j+1}^{-1}\cdots T_i^{-1}r_i^E.
 \notag
\end{align}
The first sum is empty when $j=0$.  Both sums converge in the indicated
weighted norm.  More generally, for $1<\beta<2$,
\begin{equation}
\label{geo-flat-green-bound}
 \lVert G\rVert\leq\Gamma_\beta,
 \qquad
 \Gamma_\beta=\max\left\{
 \frac{e^\beta}{1-e^{\beta-2}},\,
 \frac{e}{1-e^{1-\beta}}\right\}.
\end{equation}
For the first sum, setting $k=j-1-i$ bounds its weighted norm by
$e^\beta\sum_{k\geq0}e^{(\beta-2)k}\lVert r\rVert_\beta$.
For the second, Equation~\eqref{geo-flat-inverse-products} gives the
bound $e\sum_{k\geq0}e^{(1-\beta)k}\lVert r\rVert_\beta$.

Prescribe a small initial transverse coordinate $\xi\in\mathbb R$ and
put $\mathcal H(\xi)_j=(c^j\xi,0)$.  The desired orbit sequence is a fixed point of
\begin{equation}
\label{geo-flat-fixed-point}
 z=\mathcal H(\xi)+G\mathcal R(z),\qquad
 \mathcal R(z)_j=R_j(z_j).
\end{equation}
Choose $0<\delta<\rho$ so small that
\begin{equation}
\label{geo-flat-parameter-choice}
 \vartheta_0:=\Gamma_\beta M_0\delta\leq\tfrac12,
 \qquad \lvert\xi\rvert\leq\tfrac12\delta.
\end{equation}
On the closed ball $\lVert z\rVert_\beta\leq\delta$,
Equation~\eqref{geo-flat-remainders} yields
\[
 \lVert\mathcal R(z)\rVert_\beta\leq\tfrac12 M_0\delta^2,
 \qquad
 \lVert\mathcal R(z)-\mathcal R(\widetilde z)\rVert_\beta
 \leq M_0\delta\lVert z-\widetilde z\rVert_\beta.
\]
Consequently the right side of Equation~\eqref{geo-flat-fixed-point}
maps that ball into its ball of radius $3\delta/4$ and contracts
distances by at most $\vartheta_0$.  The contraction principle gives a unique
fixed point there, with
\begin{equation}
\label{geo-flat-sequence-decay}
 \lVert z(\xi)\rVert_\beta
 \leq\frac{\lvert\xi\rvert}{1-\vartheta_0}\leq2\lvert\xi\rvert.
\end{equation}
Substituting the sums in Equation~\eqref{geo-flat-green-sums} shows that
this fixed point satisfies Equation~\eqref{geo-flat-backward-maps} at
every step, with $s_0=\xi$.  It therefore represents an actual backward
orbit, rather than a sequence of approximate orbit points.

\emph{Dependence on the initial coordinate.}
The same contraction estimate gives
\[
 \lVert z(\xi+h)-z(\xi)\rVert_\beta
 \leq\frac{\lvert h\rvert}{1-\vartheta_0}.
\]
The map $\mathcal R$ is continuously differentiable on this sequence
ball: its derivative acts pointwise by $DR_j(z_j)$, and the uniform
second-derivative bound in Equation~\eqref{geo-flat-remainders} bounds
the remainder in the weighted norm.  In particular
$\lVert G D\mathcal R(z(\xi))\rVert\leq \vartheta_0$.  Inverting by a
Neumann series gives
\begin{equation}
\label{geo-flat-parameter-derivative}
 Dz(\xi)h
 =\bigl(I-GD\mathcal R(z(\xi))\bigr)^{-1}\mathcal H(h).
\end{equation}
For clarity, this is a genuine Fr\'echet derivative: Taylor's estimate
and the preceding Lipschitz bound give
\[
 \bigl\lVert z(\xi+h)-z(\xi)-Dz(\xi)h\bigr\rVert_\beta
 \leq\frac{\Gamma_\beta M_0}{2(1-\vartheta_0)^3}\lvert h\rvert^2.
\]
The derivative is continuous by the same uniform estimates.
At $\xi=0$, the fixed point is zero and $D\mathcal R(0)=0$, so
$Dz(0)h=((c^jh,0))_{j\geq0}$.

Define $\gamma_p(\xi)=\Psi_0(z_0(\xi))$.  Its initial derivative is
$u_0$, hence $\alpha_p(\gamma_p'(0))=1$.  Since $s_0=\xi$, it is an
embedded graph in this chart on a smaller interval.  Uniform comparison
between chart and ambient distances, together with
Equation~\eqref{geo-flat-sequence-decay}, proves
Equation~\eqref{geo-flat-arc-conclusion} at integer times.  The smooth
flow has uniformly bounded derivatives near $\mathcal K$ for times
in $[-1,0]$; decreasing the chart radius if necessary gives the same
estimate for all intervening times.  Compactness makes all constants
uniform in $p$.  Shrinking $\epsilon$ then keeps the whole backward arc
in any prescribed neighborhood of $\mathcal K$.
\end{proof}

\begin{lemma}[Flat boundary null leaves cannot occur]
\label{geo-flat-exclusion}
Let $(M,g)$ be a smooth three-dimensional Riemannian manifold and let
$N\in C^\infty(M)$ be nonnegative, with compact zero set
$Z=\{N=0\}$.  Put $P=\{N>0\}$, with boundary taken in $M$.
Suppose a neighborhood $U$ of $Z$ carries a smooth nonvanishing one-form
$\alpha$ defining a cooriented foliation by two-dimensional leaves.
Assume that every sufficiently small foliated chart can be written as
\begin{equation}
\label{geo-flat-factorization}
 \alpha=B\,\mathrm dy,\qquad B>0,\qquad
 N=B\,a(y),\qquad a\geq0,
\end{equation}
with smooth coefficients.  Let $q$ be the induced metric on a plaque
$\Sigma=\{y=y_0\}\subset\partial P$, let $K_\Sigma$ be its Gaussian
curvature, and put $b=\sqrt B$.  Suppose the following leaf equations
hold on every such plaque:
\begin{equation}
\label{geo-flat-leaf-equations}
 \Hess_\Sigma b=\tfrac12 b(K_\Sigma+3)q,
 \qquad
 \Delta_\Sigma(b^2)=6b^2-a''(y_0).
\end{equation}
Then $\Hess_gN(p)\ne0$ for every $p\in\partial P$.
Equivalently, every boundary zero in
Equation~\eqref{geo-flat-factorization} satisfies $a''(y_0)>0$.
\end{lemma}

\begin{proof}
\emph{The intrinsic vector field and the compact flat set.}
On a foliated chart define
\[
 f=\log b,\qquad Y=\nabla_\Sigma f,\qquad L=\lvert Y\rvert_g^2.
\]
On an overlap of cooriented charts the transverse coordinates satisfy
$\widetilde y=\vartheta(y)$ with $\vartheta'>0$ after restriction to
an overlap plaque.  Consequently
\[
 \widetilde B=\frac{B}{\vartheta'(y)},\qquad
 \log\sqrt{\widetilde B}
 =\log\sqrt B-\tfrac12\log\vartheta'(y).
\]
The last term has zero leaf derivative.  Thus $Y$ and $L$ are globally
defined smooth fields on $U$, even though $f$ need not be a global
function.  The vector field $Y$ is tangent to the foliation.

At a zero $y_0$ of $a$, nonnegativity gives $a'(y_0)=0$, and direct
differentiation of Equation~\eqref{geo-flat-factorization} gives
\begin{equation}
\label{geo-flat-hessian-zero}
 \Hess_gN=B\,a''(y_0)\,\mathrm dy\otimes\mathrm dy
 \quad\hbox{on the plaque }\{y=y_0\}.
\end{equation}
In particular flatness, meaning $a''(y_0)=0$, is independent of the
chosen transverse coordinate.  The boundary property is also constant
along a plaque: $N>0$ on that chart exactly where $a(y)>0$.
The set
\begin{equation}
\label{geo-flat-compact-flat-set}
 F=\partial P\cap\{\Hess_gN=0\}
\end{equation}
is therefore closed in $Z$, compact, and locally a union of whole
plaques.  It is preserved by the local flow of $Y$.

Choose a neighborhood $U_0$ of $Z$ with compact closure in $U$ and
multiply $Y$ by a smooth compactly supported cutoff equal to one near
$\overline{U_0}$.  Extend the resulting field by zero to $M$ and denote
its complete flow by $\phi_t$.  Compact support ensures completeness.
On $F$ this flow agrees with the original leafwise flow for all times:
local plaque invariance and closedness of $F$ show that
$\{t:\phi_t(p)\in F\}$ is both open and closed in $\mathbb R$ for
every $p\in F$.  In all subsequent uses outside $F$ we keep the flow
inside $U_0$, where the cutoff equals one.

\emph{The logistic identity.}
On a flat plaque, tracing the first equation in
Equation~\eqref{geo-flat-leaf-equations} gives
$\Delta_\Sigma b=b(K_\Sigma+3)$.  The second then gives
\[
 6b^2=\Delta_\Sigma(b^2)
 =2b^2(K_\Sigma+3)+2\lvert\mathrm d_\Sigma b\rvert_q^2.
\]
It follows that
\begin{equation}
\label{geo-flat-logistic-identities}
 K_\Sigma=-L,\qquad
 \Hess_\Sigma f=\frac{3-L}{2}q-
 \mathrm d_\Sigma f\otimes\mathrm d_\Sigma f,
 \qquad Y(L)=3L(1-L).
\end{equation}
For the last identity, use
$Y(L)=2\Hess_\Sigma f(Y,Y)$.

Since $F$ is compact and invariant, $L$ is bounded on every complete
orbit in $F$.  A logistic solution with $L(0)>1$ has the explicit form
\[
 L(t)=\frac{L(0)e^{3t}}{1-L(0)+L(0)e^{3t}}
\]
and blows up in finite negative time.  Thus $0\leq L\leq1$ on $F$.
Suppose for contradiction that $F$ is nonempty.  On any flat plaque
$L$ cannot vanish identically: that would give $Y=0$ and
$\Hess_\Sigma f=0$, whereas
Equation~\eqref{geo-flat-logistic-identities} would give
$\Hess_\Sigma f=\tfrac32 q$.  Hence some point of $F$ has $L>0$.
The forward logistic orbit through this point has $L(t)\to1$.
Compactness supplies an accumulation point with $L=1$, so
\begin{equation}
\label{geo-flat-unit-set}
 F_1=F\cap\{L=1\}
\end{equation}
is a nonempty compact set invariant under the full flow.

\emph{Leaf and transverse derivative rates.}
Along $F_1$, Equation~\eqref{geo-flat-logistic-identities} becomes
\begin{equation}
\label{geo-flat-unit-hessian}
 \Hess_\Sigma f=q-\mathrm d_\Sigma f\otimes\mathrm d_\Sigma f,
 \qquad \lvert Y\rvert_g=1.
\end{equation}
For a leaf-tangent vector $v(t)=D\phi_t(p)v(0)$, $p\in F_1$, the
usual derivative of its squared length within the leaf is
\[
 \frac{\mathrm d}{\mathrm dt}\lvert v(t)\rvert_g^2
 =2\Hess_\Sigma f(v(t),v(t))
 =2\bigl(\lvert v(t)\rvert_g^2-
             \langle Y,v(t)\rangle_g^2\bigr).
\]
This lies between zero and $2\lvert v(t)\rvert_g^2$.  Integration gives
the two tangent inequalities in
Equation~\eqref{geo-flat-rate-hypotheses}.
Furthermore, throughout $U_0$, tangency of $Y$ implies $Y(y)=0$ and
\begin{equation}
\label{geo-flat-alpha-and-n-evolution}
 \mathcal L_Y\alpha=Y(B)\,\mathrm dy=2L\alpha,
 \qquad Y(N)=2LN.
\end{equation}
Here $Y(\log B)=2Y(f)=2L$; the second identity also holds where $N=0$,
by the smooth factorization.  Since $L=1$ along $F_1$, the first
identity gives
$\alpha_{\phi_t(p)}\circ D\phi_t(p)=e^{2t}\alpha_p$.
All hypotheses of Lemma~\ref{geo-flat-fast-arc} therefore hold with
$\mathcal K=F_1$ and $E=\ker\alpha$.

\emph{The two incompatible decay bounds.}
Fix $p\in F_1$ and apply Lemma~\ref{geo-flat-fast-arc}.  By compactness
and continuity choose a neighborhood $V$ of $F_1$, with closure in
$U_0$, on which $L\leq5/4$.  Shrink the resulting transverse arc so
that all its backward trajectories remain in $V$.  If one of its
points $q=\gamma_p(\xi)$ had $N(q)>0$, the second equation in
Equation~\eqref{geo-flat-alpha-and-n-evolution} would give, for every
$t\geq0$,
\begin{equation}
\label{geo-flat-n-lower-bound}
 N(\phi_{-t}(q))
 =N(q)\exp\left(-\int_0^t2L(\phi_{-s}(q))\,\mathrm ds\right)
 \geq N(q)e^{-5t/2}.
\end{equation}
In particular the trajectory remains positive at each finite time.

On the other hand, $N$ and $\mathrm dN$ vanish at every point of $F_1$,
because $N$ is smooth and nonnegative.  Its Hessian is uniformly bounded
on a compact neighborhood of $F_1$.  Taylor's Theorem in uniform
normal-coordinate balls therefore gives
\[
 0\leq N(x)\leq C_0d_g(x,z)^2
 \quad\text{when }z\in F_1\text{ and }x\text{ is sufficiently close to }z.
\]
Apply this with $z=\phi_{-j}(p)$, $x=\phi_{-j}(q)$ and use
Equation~\eqref{geo-flat-arc-conclusion}.  Uniformly for all integers
$j\geq0$,
\begin{equation}
\label{geo-flat-n-upper-bound}
 N(\phi_{-j}(q))\leq C_1\lvert\xi\rvert^2e^{-3j}.
\end{equation}
Equations~\eqref{geo-flat-n-lower-bound} and
\eqref{geo-flat-n-upper-bound} imply
$N(q)\leq C_1\lvert\xi\rvert^2e^{-j/2}$ for every $j$, a
contradiction.  Thus $N$ vanishes on the whole small arc.

In a foliated chart at $p$, transversality gives
$\frac{\mathrm d}{\mathrm d\xi}y(\gamma_p(\xi))|_{\xi=0}\ne0$.
After restriction, the transverse coordinate of the arc covers an open
interval containing $y(p)$.  Equation~\eqref{geo-flat-factorization}
and positivity of $B$ then imply that $a$ vanishes on that interval.
This gives an open neighborhood of $p$ on which $N=0$, contradicting
$p\in\partial P$.  Hence $F$ is empty.  Finally $a''\geq0$ at any zero
of the nonnegative function $a$, so
Equation~\eqref{geo-flat-hessian-zero} proves the stated strict
positivity at every boundary zero.
\end{proof}


\begin{proposition}[No interior null set]
\label{geo-no-null-set}
One has $Z=\varnothing$, and hence $P=M$.
\end{proposition}
\begin{proof}
Apply Lemma~\ref{geo-flat-exclusion} with $\alpha=X^\flat$; its hypotheses
are precisely the compactness in Lemma~\ref{geo-positive-collars} and the
factorization and equations in Lemma~\ref{geo-null-leaf-equations}.
Every boundary zero therefore satisfies $a''>0$.
It is therefore an isolated zero in the transverse coordinate, and its
neighborhood consists of a smooth two-sided zero leaf with $P$ on both
sides.  There can be no open null region.  Indeed, if the interior of $Z$
were nonempty, it would have a boundary in the connected manifold $M$,
since $P$ is nonempty.  Such a boundary point lies in $\partial P$ and would
have the just-described neighborhood with no open zero region, a
contradiction.  Thus all of $Z$ equals $\partial P$, and it is a compact
embedded hypersurface.  The local isolation and compactness give only
finitely many components.  The nonvanishing global form $X^\flat$ coorients
each of them.

On a component of $Z$, the local positive functions $B_1/a''$ patch to a
global smooth function.  To check this, another oriented transverse
coordinate has the form $\widetilde y=\phi(y)$ with $\phi'>0$.  At a zero
leaf,
\[
 \widetilde B_1=B_1/\phi',\qquad
 \widetilde a=\phi'a,\qquad
 \widetilde a''=a''/\phi',
\]
where the last derivatives use $\widetilde y$ and $a=a'=0$.
The quotient is therefore invariant.  Since $a''$ is constant along each
plaque, Equation~\eqref{geo-boundary-leaf-system} gives
\[
 \Delta_{q_y}\left(\frac{B_1}{a''}\right)
                 =6\frac{B_1}{a''}-1.
\]
The maximum and minimum principles on the compact component imply
$B_1/a''=1/6$.  Its leaf derivative is zero; thus $b_1$ is constant along
each plaque, and the first equation in
Equation~\eqref{geo-boundary-leaf-system} gives $\mathcal K_y=-3$.

Cut the original manifold along these finitely many two-sided components
and truncate its ordinary end.  The resulting compact orientable
three-manifold $L$ is connected and has interior homotopy equivalent to
$P$: the end and each of the original and cut boundary collars are product
collars, which may be shortened without changing homotopy type.
Proposition~\ref{geo-simple-connectivity} therefore gives
$H_1(L;\mathbb Z)=H^1(L;\mathbb Z)=0$.  Poincar\'e--Lefschetz duality and
the boundary long exact sequence give
\[
 H_2(L,\partial L;\mathbb Z)\cong H^1(L;\mathbb Z)=0,
 \qquad
 0=H_2(L,\partial L;\mathbb Z)\longrightarrow
 H_1(\partial L;\mathbb Z)\longrightarrow H_1(L;\mathbb Z)=0.
\]
Consequently every component of the orientable boundary is a sphere.  In
particular each cut copy of a component of $Z$ is a sphere, incompatible
with its constant Gaussian curvature $-3$ by Gauss--Bonnet.  Thus $Z$ must
have been empty.
\end{proof}

\subsection{The regular boundary of the quotient}

\begin{lemma}[Static boundary collars and the time primitive]
\label{geo-boundary-collars}
\label{geo-static-collar}
The metric $h$ has the regular completion at $S$ stated in
Equation~\eqref{geo-quotient-boundary}.  If $u|_S=0$, its boundary structure
is the original smooth boundary structure and $\mathcal A$ extends smoothly
to $S$.  If $u|_S>0$, use $N$ as the original defining function and use
$\lambda=\sqrt N$ as the quotient defining function.  The quotient metric
extends smoothly by reflection $\lambda\mapsto-\lambda$; its diagonal
blocks are even and mixed block odd.  In the original collar,
\begin{equation}
 \mathcal A=\frac1{2\kappa_h}\dd\log N+\beta,
 \qquad \beta\text{ is smooth and closed up to }S.
 \label{geo-positive-A-collar}
\end{equation}
In quotient normal distance $\rho$ from $S$, both cases satisfy
\begin{equation}
 h=\dd\rho^2+q_0+O(\rho^2),\qquad
 \lambda=\kappa_h\rho+O(\rho^3),
 \label{geo-normal-boundary-jets}
\end{equation}
with smooth one-sided coefficients.
For a global primitive $\dd H=-\mathcal A$, one has, respectively,
\begin{equation}
 H\text{ smooth up to }S,
 \qquad\text{or}\qquad
 H=-\frac1{2\kappa_h}\log N+B_0,
 \quad B_0\text{ smooth up to }S.
 \label{geo-time-collar}
\end{equation}
\end{lemma}
\begin{proof}
In the zero-lapse case put $d=a^2-s^2|W|_g^2$ in
Equation~\eqref{geo-zero-boundary-jet}.  Then
\[
 h=g+s^2d^{-1}W^\flat\otimes W^\flat,
 \qquad \lambda=s\sqrt d,
 \qquad \mathcal A=d^{-1}W^\flat,
 \qquad d|_S=\kappa_h^2.
\]
These are smooth in the original collar and give
$h|_{TS}=g|_{TS}$ and $\partial_{\nu_h}\lambda=\kappa_h$.
The static Hessian equation extends to $S$ and gives $\Hess_h\lambda=0$
there.  Its tangential restriction is $\kappa_h\mathrm{II}_h$, proving
$\mathrm{II}_h=0$.

If $u|_S>0$, Equation~\eqref{geo-positive-boundary-jet} makes $N$ a smooth
original defining function.  In coordinates $(N,y^1,y^2)$, smooth division
of the tangential components gives
\[
 X^\flat=a_0(N,y)\dd N+N\beta_A(N,y)\dd y^A,
 \qquad a_0(0,y)=\frac1{2\kappa_h}.
\]
The change $N=\lambda^2$ gives the metric coefficients
\begin{align}
 h_{\lambda\lambda}&=4\lambda^2g_{NN}+4a_0^2,
 &h_{\lambda A}&=2\lambda(g_{NA}+a_0\beta_A),
 &h_{AB}&=g_{AB}+\lambda^2\beta_A\beta_B,
 \label{geo-even-odd-collar}
\end{align}
with the coefficients on the right evaluated at $(\lambda^2,y)$.  They are
smooth across $\lambda=0$ with the stated parity.  At zero the normal block
is $\kappa_h^{-2}$, the mixed block vanishes, and the tangential block is
$g|_{TS}$, so this is a nondegenerate smooth Riemannian extension.
Reflection is an isometry; its fixed hypersurface is totally geodesic.
The normal derivative of $\lambda$ is $\kappa_h$.
In either case, quotient Gaussian coordinates and $\mathrm{II}_h=0$ give
$\partial_\rho q_\rho|_0=0$.  The static equation gives
$\partial_\rho^2\lambda|_0=0$, proving
Equation~\eqref{geo-normal-boundary-jets}.

Since $a_0-(2\kappa_h)^{-1}$ vanishes at $N=0$, division by $N$ also gives
Equation~\eqref{geo-positive-A-collar}.  Its remainder is closed by
$\dd\mathcal A=0$ in the interior and smoothness at the boundary.
Simple connectivity and Lemma~\ref{geo-vanishing-twist} give a global $H$
with $\dd H=-\mathcal A$.  In either case subtract the displayed singular
term, if present.  The remaining differential is smooth in the original
collar.  Fixing values on one interior collar section and integrating the
normal derivative to $S$ gives a smooth extension; its tangential
derivatives agree with the prescribed differential by continuity.  This
proves Equation~\eqref{geo-time-collar}.
\end{proof}

\begin{proof}[Proof of Theorem~\ref{geo-static-quotient}]
Propositions~\ref{geo-simple-connectivity} and \ref{geo-no-null-set} give
$P=M$ and its simple connectivity.  Lemmas~\ref{geo-vanishing-twist},
\ref{geo-static-on-P}, and \ref{geo-boundary-collars} prove the stated
closedness, equations, and boundary data.

The quotient completion has the same underlying topology as $\Omega$:
in the positive-lapse case $N\mapsto\sqrt N$ changes the boundary smooth
structure but is a homeomorphism of the closed collar, and its restriction
to $S$ is the identity.  To check completeness, let a path have finite
$h$-length.  Since $h\geq g$ in the interior, it is a $g$-Cauchy path and
has a limit in the complete original manifold with boundary.  An interior
limit is an ordinary smooth quotient point because $N>0$ there.  A limit
on $S$ is included by Lemma~\ref{geo-boundary-collars}; in its smooth
quotient collar the same path tends to that boundary point.  The unique
ordinary end has infinite $g$-distance, and hence infinite $h$-distance.
These exhaust all possible finite-length escapes and prove completeness.

Finally, Equation~\eqref{geo-input-end} gives
$N=V_0^2(1+O_2(r^{-\tau'}))$.  Since spatial lowering by $g$ preserves the
stated end orders,
\[
 N^{-1}X^\flat\otimes X^\flat=O_2(r^{-2\tau'}),
 \qquad \sqrt N-V_0=O_2(r^{1-\tau'}).
\]
These are Equation~\eqref{geo-quotient-end}.  In particular the attached
boundary is compact, the quotient has the single designated end, and no
additional completion points have been introduced.
\end{proof}
